% Generated draft: requires mathematical and source review.
\documentclass[11pt]{article}
\usepackage[margin=25mm]{geometry}
\usepackage{fontspec}
\setmainfont{Latin Modern Roman}
\usepackage{amsmath,amssymb,mathtools,unicode-math}
\setmathfont{Latin Modern Math}
\usepackage{enumitem,needspace,xurl,hyperref}
\setlength{\parindent}{0pt}
\setlength{\parskip}{5pt}
\setlength{\emergencystretch}{3em}
\newcommand{\N}{\mathbb{N}}
\newcommand{\Z}{\mathbb{Z}}
\newcommand{\Q}{\mathbb{Q}}
\newcommand{\R}{\mathbb{R}}
\newcommand{\C}{\mathbb{C}}
\newcommand{\F}{\mathbb{F}}
\newcommand{\A}{\mathbb{A}}
\newcommand{\PP}{\mathbb{P}}
\newcommand{\E}{\mathbb{E}}
\newcommand{\eps}{\varepsilon}
\newcommand{\dd}{\,\mathrm{d}}
\begin{document}
\section*{UnsolvedMath ID: 701643}
\textbf{Source problem code:} catalogue entry
\subsection{Definitions and mathematical statement}
Let \(n\geq 3\) be an integer, and let \(G_\Gamma\) be a partially commutative soluble group of derived length \(n\), with defining graph \(\Gamma\), in the sense intended by the source. Write \(V(\Gamma)\) for its vertex set and \(x_v\in G_\Gamma\) for the distinguished generator corresponding to \(v\in V(\Gamma)\). Define the derived series by
\[
G_\Gamma^{(0)}=G_\Gamma,\qquad G_\Gamma^{(k+1)}=[G_\Gamma^{(k)},G_\Gamma^{(k)}].
\]
Derived length \(n\) means that \(G_\Gamma^{(n)}=\{1\}\) and \(G_\Gamma^{(n-1)}\neq\{1\}\). For each vertex \(v\), put
\[
N_\Gamma(v)=\{u\in V(\Gamma):u\text{ is adjacent to }v\},\qquad
C_{G_\Gamma}(x_v)=\{g\in G_\Gamma:gx_v=x_vg\}.
\]
The source asks whether, for every such integer \(n\), every such group with its defining graph, and every \(v\in V(\Gamma)\),
\[
C_{G_\Gamma}(x_v)=\bigl\langle x_u:u\in N_\Gamma(v)\bigr\rangle.
\]
\emph{Unresolved qualifications.} The source refers to Problem 17.104 for the analogous definition when \(n=2\), but that definition is not supplied. Thus the precise class of groups quantified over is not fully specified here. The adjacency convention is also unspecified. Since \(x_v\in C_{G_\Gamma}(x_v)\), it matters whether the intended generating set includes \(x_v\) itself. The displayed question preserves the supplied wording; replacing its right-hand side by \(\langle x_v,x_u:u\in N_\Gamma(v)\rangle\) would be a different formulation requiring source clarification.
\subsection{Short English statement}
For a partially commutative soluble group of derived length at least three, must the centralizer of each vertex generator be generated by the generators at adjacent vertices? The intended group definition and whether the vertex itself belongs to this generating set require clarification.
\subsection{Sources}
\begin{itemize}
\item Supplied local UnsolvedMath registry record: permanent ID \(701643\), source problem code \texttt{catalogue entry}, title ``Kourovka Notebook Problem 19.95.'' Upstream revision: \texttt{d371ee917ac9647fbb06cd6fc529deec6f3dda82}. The registry labels the problem open; this is supplied metadata, not an independently verified status assessment.
\item \emph{The Kourovka Notebook (Unsolved Problems in Group Theory)}, Problem 19.95, proposed by E. I. Timoshenko in 2018. The supplied background identifies the 19th issue (2018) and the 21st issue edition (2026), with identifier \texttt{arXiv:1401.0300v48}. The supplied collection description names E. I. Khukhro and V. D. Mazurov as editors.
\item The statement refers to Kourovka Notebook Problem 17.104 for the analogous definition in derived length two. Its text is not included in the supplied record. No source was independently consulted for this draft.
\end{itemize}
\end{document}
